% Include as a chapter or section fragment. All matrices are printed in the JSON.
\section{An explicit fixed family from the six-root identity}\label{sec:six}

This construction represents the complete reciprocal C1 parent on an
explicit source domain and an independent testing rectangle. It proves
membership in a particular fixed Gram cone on that domain. It makes no
assertion that the domain covers all C1 sources or any of the omitted C2
source families.

\paragraph{Source notation.}
Use (P1), (SP3), $Q=Q_b=1-b$, and the source quantities
\[
 L=\lambda/\sigma,\qquad s=\sqrt{1+L^2},\qquad d_G=4D/\sigma^2.
\]
Within this section $s$ is this pure-source normalizer, not the
testing square root used in the preceding section.


\paragraph{The exact full-state Gram.}
For independent testing coordinates $u,w,W,X,Y,Z$ set
\[
\begin{aligned}
 N_x&=X+uhZ,&N_y&=(Y-bwZ)/q,&R&=wZ/q-uqW,\\
 \Gamma&=\sigma Z-uwW,& A_\theta&=q(u^2-x)W,
 &\widehat V&=suw(u^2-x)W.
\end{aligned}
\]
Let $v=(N_x,N_y,R,\Gamma,A_\theta,\widehat V)^T$.
The exact identity used here is
\[
 F_w=E_{\rm ret}+v^T\!\left(G_0+
                  \sum_{\gamma\in\mathcal A}u^{\gamma_1}w^{\gamma_2}E_\gamma
                          \right)v,\qquad
 \mathcal A=\{00,10,01,11,20,02,21,30\}.
 \tag{SR1}
\]
Here
\[
G_0=\begin{pmatrix}
2&0&0&0&0&0\\
0&2&2b-1&0&0&0\\
0&2b-1&b^2+Q^2&0&0&0\\
0&0&0&d_G&0&L/(2s)\\
0&0&0&0&4\kappa H&0\\
0&0&0&L/(2s)&0&1
\end{pmatrix}.
\]
All unspecified entries of the symmetric matrices below are zero.
Indices follow the six roots just printed, starting at 1.
\[
\begin{array}{c|c|l}
\gamma&(i,j)&(E_\gamma)_{ij}\\\hline
00&(3,3)&z[8bD+8bp\sigma+4b\sigma(2C+2D+1)(r-1)]\\
00&(2,3)&2z[\sigma(2C+1)(r-1)+2D(1+\sigma(r-1))+2p\sigma]\\
00&(4,6)&-sx/2\\\hline
10&(1,4)&-2\nu\\
10&(2,5)&A_y/2\\
10&(3,5)&(\chi+\lambda+s^2\sigma x)/2\\\hline
01&(2,4)&qR_0/2\\
01&(3,4)&bqR_0/2\\
01&(1,5)&qB_x/2\\\hline
11&(1,2)&-2q(2m+2D+1)\\
11&(1,3)&qR_x/2\\\hline
20&(2,2)&2z(1+2m)\\
20&(3,3)&zk_1\\
20&(2,3)&2z[2Dh+2p\sigma+(1+2m)g]\\
20&(4,4)&L\\
20&(5,5)&-s^2\sigma\\
20&(4,6)&s/2\\\hline
02&(2,2)&2z(1+2D)\\
02&(3,3)&2zb^2(2D+1)\\
02&(2,3)&2zb(2D+1)\\\hline
21&(2,4)&-4pq\\
21&(3,4)&-4bpq\\\hline
30&(3,5)&-s^2\sigma/2
\end{array}
\tag{SR2}
\]
In particular, the $N_x\Gamma$, $N_y\Gamma$, and $R\Gamma$
entries retain both complete state-normal couplings. No determinant
substitution has replaced the four-state parent.

For completeness, define $X_e=X+ewW$,
$J_C=Y/q-buqW$, and $J_\ell=(Y-wZ)/q+QuqW$. The retained form is
\[
\begin{aligned}
E_{\rm ret}={}&4\beta N_x^2+(4\beta+2+4m)(wN_x)^2
 +2(2D+1)(uX_e)^2+4bpz(u\Gamma)^2\\
&+[4a+2z(2C+1)-1]J_C^2+(4\ell-1)J_\ell^2\\
&+4a(uJ_C)^2+(4a+4Cz)(wJ_C)^2
       +4a(uwJ_C)^2+4\ell(uJ_\ell)^2.
\end{aligned}\tag{SR3}
\]
All ten coefficients are nonnegative when $4a,4\ell\ge1$.

To prove (SR1), start with the full-parent identity SP6. With
$V=uw(u^2-x)W$, direct multiplication gives
\[
\begin{aligned}
 K_AA_\theta^2+\mathscr S
 &=4\kappa H A_\theta^2+Lu^2\Gamma^2+L\Gamma V
                                       +\lambda uA_\theta R,\\
0&=V^2-\sigma u^2A_\theta^2-u\sigma(u^2-x)A_\theta R
                                      +(u^2-x)V\Gamma.
\end{aligned}
\]
Add $s^2=1+L^2$ times the second identity, replace $V$ by
$\widehat V/s$, and collect the eight testing monomials. Clipping
the two radial squares at one gives (SR3). This proves (SR1)--(SR3)
algebraically at either sign of $\lambda$.

\paragraph{Explicit source guards.}
For a real symmetric matrix set
$\|T\|_{\rm row}=\max_i\sum_j|T_{ij}|$. Choose the fixed rational
constants
\[
 \delta=2^{-12}=1/4096,\qquad \eta=1/56.
\]
The source admission conditions, in addition to the native domain, are
\[
 4a\ge1,\quad4\ell\ge1,\quad d_G\ge1,\quad4\kappa H\ge1,
 \tag{SR4}
\]
\[
 \|E_{00}\|_{\rm row}\le1/12,\qquad
 \delta^{|\gamma|}\|E_\gamma\|_{\rm row}\le1/56
                   \quad(\gamma\in\mathcal A\setminus\{00\}).
 \tag{SR5}
\]
These are finite, explicit inequalities in source-only rational and
algebraic functions; no testing point is used to select them.

The radial block has determinant one and trace at most three, so it
is at least $I_2/3$. An explicit pure-source proof is
\[
 (N_y+bR)^2+(N_y-QR)^2
 =\frac{N_y^2+R^2}{3}
 +\frac{5}{3}\left(N_y+\frac{3(2b-1)}{5}R\right)^2
 +\frac{1+6bQ}{15}R^2.
\]
The $(\Gamma,\widehat V)$ block is at least $I_2/2$: its diagonals
are at least one and $|L/s|\le1$. Thus $G_0\succeq I_6/3$.
Consequently
\[
 B=G_0+E_{00}-I_6/8\succeq I_6/8.
\]
For the seven nonconstant monomials put
\[
 y_\gamma=(u/\delta)^{\gamma_1}(w/\delta)^{\gamma_2},\qquad
 T_\gamma=\delta^{|\gamma|}E_\gamma.
\]
On the fixed independent testing box $0\le u,w\le\delta$ one has
$0\le y_\gamma\le1$. The exact matrix identity is
\[
G_0+\sum_\gamma u^{\gamma_1}w^{\gamma_2}E_\gamma
=B+\frac12\sum_{\gamma\ne00}
 \{(\eta I_6+T_\gamma)(1+y_\gamma)
       +(\eta I_6-T_\gamma)(1-y_\gamma)\}.
\tag{SR6}
\]
The scalar equality $7\eta=1/8$ proves the constant coefficient.
The source matrices $\eta I_6\pm T_\gamma$ are positive semidefinite
by (SR5). All factors $1\pm y_\gamma$ are nonnegative in the testing
field. Factor each source matrix over its own real closed field.
This proves a separated tensor SOS with the \emph{fixed} fourteen
testing radicals $\sqrt{1\pm y_\gamma}$. No adaptive constant $M$
or source-dependent testing function occurs.

\paragraph{A full uniform source rectangle for the nonpositive branch.}
The guards hold for every native source satisfying
\[
 0<x,\sigma\le\varepsilon:=2^{-10}=1/1024,\qquad
 0<z\le1,\qquad \lambda\le0.
\tag{SR7}
\]
No inverse of $\beta$ occurs, so its literal zero face is included.
The proof below forces $D>0$ on this nonpositive branch; the $D=0$
source face is therefore assigned to the complementary positive
branch, rather than being cancelled or omitted.

Here are elementary bounds establishing the claim. From the
definition of $p$, $P\ge1$ and $0<E<1$, one has $p\le1+\varepsilon<2$.
Since every term preceding $-4bp\sigma z$ in $\lambda$ is nonnegative,
\[
 H\le2,\quad h\le2\sigma,\quad e\le2x,\quad
 \beta\theta^2\le2\sigma^2,\quad -8\le L\le0,\quad1\le s\le9.
\]
In particular $\theta\ge1/2$ and $\beta\theta\le4\sigma^2$.
The source formulas give
\[
 1\le P\le1+4\varepsilon,\quad
 1-\varepsilon\le M\le1+\varepsilon,\quad
 1-\varepsilon\le J\le1+\varepsilon,\quad
 B_s=P+xz(J-\sigma)\le1+6\varepsilon.
\]
The exact identity $2\theta P\beta=hB_s-\sigma H_0$, where
$H_0=1+zQ\ge1$, implies $H\ge3/4$. Also
\[
 3/8\le\kappa\le1,\quad1/4\le p\le2,\quad
 1/4\le l\le1,\quad C(1-r)\le1,\quad m\le1/8.
\]
For the lower bound on $l$, use
$2Pl\ge M(1-e)\ge(1-\varepsilon)(1-2\varepsilon)$.
Its upper bound follows from
$eR_s(1+\sigma z)+M(1-e)\le(1+\varepsilon)^2$.
These estimates already prove the first, second and fourth guards in
(SR4).

To bound $D$, use the exact BE4 identity
\[
 2P(JD+x\theta\beta)=\sigma^2\mathcal H,
\]
where
\[
\mathcal H=R_sH_0+z\{(1-r)xh(2+xr)
       -e[2+x(1+h+r\theta)]+Ke^2\}.
\]
The preceding inequalities give $1-5\varepsilon\le\mathcal H\le3$,
and hence $D\le3\sigma^2$. Moreover
\[
 d_G\ge
 \frac{2(1-5\varepsilon)/(1+4\varepsilon)-16\varepsilon}
      {1+\varepsilon}\ge1;
\]
the last inequality follows from
$1-31\varepsilon-68\varepsilon^2\ge0$. Thus (SR4) is proved.

The remaining coefficient estimates are deliberately loose:
\[
 |\nu|\le3,\quad |R_0|\le29,\quad |A_y|\le22,
 \quad|\chi|\le24,\quad |R_x|\le6.
\]
They follow by substituting the displayed bounds in the five explicit
source formulas. For example $|g|\le1$ and
$pb+lQ\le2$; all appearances of $C$ in (SR2) outside whole
positive squares occur as $C(1-r)$, whose bound is one.
Direct row sums in (SR2) then give
\[
 \|E_{00}\|_{\rm row}\le64\varepsilon,\qquad
 \|E_\gamma\|_{\rm row}\le64\quad(\gamma\ne00).
\]
For the constant row, the sharper estimates are
$|(E_{00})_{33}|+|(E_{00})_{23}|\le
42\sigma+12D(1+\sigma)<43\varepsilon$ and
$|(E_{00})_{46}|\le9x/2$.
Thus $64\varepsilon=1/16<1/12$ and
$64\delta^{|\gamma|}\le1/64<1/56$, proving (SR5).
These are finite real inequalities, so they hold over arbitrary
real closed source and testing fields by the same algebraic proof.

\paragraph{The fixed testing span.}
Set
\[
 \phi_{10}=(W,X,Y,Z,uZ,wZ,uW,uwW,u^2W,u^3wW)^T.
\]
The exact $6\times10$ source matrix $T$ satisfying $v=T\phi_{10}$ is
\[
T=\begin{pmatrix}
0&1&0&0&h&0&0&0&0&0\\
0&0&1/q&0&0&-b/q&0&0&0&0\\
0&0&0&0&0&1/q&-q&0&0&0\\
0&0&0&\sigma&0&0&0&-1&0&0\\
-qx&0&0&0&0&0&0&0&q&0\\
0&0&0&0&0&0&0&-sx&0&s
\end{pmatrix}.
\]
The first principal block is $T^TBT$; the other fourteen principal
blocks are
\[
 \tfrac12T^T(\eta I_6\pm T_\gamma)T
 \quad(\gamma\ne00)
\]
on testing vectors $\sqrt{1\pm y_\gamma}\,\phi_{10}$.
All these matrices are positive semidefinite by (SR4)--(SR6).

The independent algebraic testing cover adds fourteen variables
$\zeta_{\gamma,\pm}$ with
\[
 \zeta_{\gamma,\pm}^2=1\pm y_\gamma,\qquad
 b_u^2=u(\delta-u),\qquad b_w^2=w(\delta-w),
\]
and the equations
\[
 W^2+X^2+Y^2+Z^2=1,\qquad
 \kappa_t^2(1+u^2)(1+w^2)=1 .
\]
These eighteen equations in twenty-three real variables define a
compact algebraic set covering the entire test rectangle and state
sphere. Multiply every testing vector by $\kappa_t$ to obtain the
normalized physical slack. Each radical multiplies a whole Gram
block, so either sign of every radical is allowed. Axes and zero
state coordinates are literal specializations. The gauges (SG)
and real/imaginary splitting cover negative angles and complex states.

\paragraph{Preserving the elementary root structure: the primary pencil.}
Keep all fifteen
principal blocks on $\phi_{10}$. Absorb the retained $N_x,J_C,J_\ell$
squares into the first principal block; these three rows already lie
in $\operatorname{span}\phi_{10}$. Group the seven remaining retained
rows only as follows:
\[
\begin{array}{c|c|c}
\text{fixed testing vector}&\text{retained roots}&\text{block order}\\\hline
(wX,uwZ)&wN_x&2\\
(uX,wY,uwW)&uX_e,\ wJ_C&3\\
(uZ,uwY,u^2wW)&u\Gamma,\ uwJ_C&3\\
(uY,uwZ,u^2W)&uJ_C,\ uJ_\ell&3
\end{array}
\]
Each block is the sum of the nonnegative source multiples of the
displayed coefficient-vector outer products. Thus every block is
positive semidefinite over the source field. This compilation has
exactly nineteen blocks, fifteen of order ten and four of orders
$(2,3,3,3)$, with total order $161$ and maximum block order ten.
It preserves the testing support of each existing root rather than
adding arbitrary cross terms among all seventeen base monomials.

The evaluated rational source pencil is
\texttt{six\_root\_structured\_source\_pencil.coo.json}.
Its visible coordinates are $(d,m,C,A,D,\alpha,\beta)$; it has
666 free auxiliary coordinates and 2311 nonzero upper-triangular
rational coefficient records. The common feature-product expansion
has 180 coefficient equations and 846 raw upper-triangular Gram
variables. Exact elimination has rank 180, leaving $846-180=666$
free variables. The matching map and every resulting source and
auxiliary coefficient are supplied and exactly checked.

Here both directions of the local cone representation are algebraic.
An affine-pencil witness supplies positive semidefinite Gram blocks;
the checked coefficient map proves their weighted sum equals the
complete parent. Conversely, any collection of Gram blocks giving
that parent is a solution of the 180 printed linear equations. Take
its entries at the printed free positions as the 666 auxiliaries.
The exact pivot equations then reconstruct every other entry, so the
printed affine pencil represents those same Gram blocks. Thus this
pencil is exactly the specified local fixed Gram cone. This is a
two-sided statement about that cone, not a quantum reverse-inclusion
claim.

An admitted rational source for the general guards is
$\sigma=10^{-4}$, $x=10^{-8}$, $b=1/2$, $z=1$, $h=3/20000$.
It has $\lambda>0$ and is therefore not an example from the
$\lambda\le0$ subdomain. The nineteen source Gram matrices at
this point are printed in the structured fixture, with $331$
nonzero rational upper-triangle entries. The source radical
$s$ cancels from their evaluated entries. This fixture is an
additional exact consistency check; the uniform admission theorem
is proved by the preceding inequalities.

The result is a fixed local testing cone with a proved admission
theorem. It is not a primal affine lift of the full eight-coordinate
set. To obtain such a lift one still needs a finite cover of every
valid source and every independent test by compatible fixed testing
families, including the remaining C1 and C2 boundaries.
